%!TeX root=./tese.tex
% Introduction. Quantitative claims quoted here are established in the body
% chapters and reproducible from the workspace's versioned scripts
% (analysis/trace_eda, predictors, collector_artifacts, simulation).
% One \emph{[Placeholder: ...]} remains (validation anchors, resolved by
% the validation chapter); all others were resolved 2026-07-08 from the
% completed baseline/closed-loop experiments (see progress/results-ledger.md).
\chapter{Introduction}
\label{cap:intro}

Large language models are served interactively, under explicit latency
targets, from fleets of GPUs whose cost dominates the economics of the
service. The systems that serve them have evolved toward decomposition.
Inference is split into stages of distinct computational character, and
each stage is provisioned separately: prefill and decode run across
dedicated worker
pools~\cite{Zhong2024DistServe,Patel2024Splitwise,Qin2025Mooncake}, above
engines that batch and schedule
tokens~\cite{Kwon2023vLLM,Agrawal2024SarathiServe}, under frameworks that
orchestrate the pools at datacenter scale~\cite{Dynamo2025}. Multimodal
models extend the decomposition by one
stage: before any text token is processed, a vision encoder must convert
the request's images into the embedding tokens the language model
consumes, and modality-disaggregated architectures give that encoder a
pool of its own~\cite{Qiu2025ModServe}. Disaggregation makes each stage
individually efficient --- and multiplies the provisioning problem. Where
a monolithic deployment asked one sizing question, a disaggregated
deployment asks one per pool, continuously, under load that moves by the
minute.

This thesis studies how the pools of a modality-disaggregated multimodal
serving system should be scaled, and defends a specific answer: each pool
must be provisioned on a forecast of \emph{its own} demand signal. The
argument rests on four results, each carried by a chapter of this
document. In a production multimodal trace, vision demand is nearly
decoupled from the request rate a conventional autoscaler watches, and the
burstiest demand signal is simultaneously the most persistent --- the most
forecastable (Chapter~\ref{cap:trace-eda}). Pretrained time-series
foundation models convert that persistence into usable forecasts,
zero-shot, beating classical baselines on every demand signal at every
resolution tested (Chapter~\ref{cap:predictors}). A cost model for the
vision stage, calibrated on silicon data collected for this thesis and
audited against physical measurement --- the collected attention operator
agrees within ${\approx}3\%$, and the composed forward is bracketed by a
fully attributed measurement band whose direction is conservative for
every claim built on it --- makes the encoder pool's simulated timing
defensible (Chapter~\ref{cap:aiconfigurator}). Composed with a calibrated
discrete-event simulation of the language pool, the two-pool system
reproduces the decoupling as measured behavior: the pools saturate at
trace scalings an order of magnitude apart on steady traffic, and the
trace's largest burst breaks them in structurally different ways --- the
language pool through a prefill floor that one static parallelism change
removes, the vision pool through queueing at 29\% utilization that only
capacity removes and no rate- or utilization-based monitor sees
(Chapter~\ref{cap:twopool}). Closing the loop, per-pool forecasts driving
worker schedules hold the tail latency objective through burst onsets
that defeat an identically informed reactive controller at every tested
cold-start lead time, at a fraction of the statically provisioned
fleet's GPU-hours (baselines and closed-loop chapters, in preparation;
results summarized in Section~\ref{sec:intro-contrib}).

\section{Context: serving multimodal models by disaggregation}
\label{sec:intro-context}

An LLM answers a request in two phases of opposite character: a
\emph{prefill} pass that processes the whole prompt at once ---
compute-bound and parallel --- and a \emph{decode} loop that generates one
token at a time, bound by memory bandwidth and exposed to the user as
inter-token latency. Service quality is specified against this anatomy,
most commonly as bounds on the time to first token (TTFT, dominated by
queueing and prefill) and the time per output token (TPOT, decode; this
document reports the equivalent per-gap measurement as inter-token
latency, ITL). Formally such bounds are service-level \emph{objectives}
(SLOs); following common serving-systems usage, this document writes
``the evaluation SLA'' for the specific set of tail bounds it grades
every policy against, and SLO where the general concept is meant.
Meeting such bounds economically on shared GPU fleets has pushed serving
architecture from replicated monoliths toward disaggregation: prefill and
decode run in separate pools sized to their own
bottlenecks~\cite{Zhong2024DistServe,Patel2024Splitwise} --- a split now
operated at production scale~\cite{Qin2025Mooncake} --- while
datacenter-scale frameworks route requests, move KV caches between pools,
and scale workers as load shifts~\cite{Dynamo2025}.

A multimodal request adds a third phase that precedes both. A vision
encoder --- in current practice a Vision
Transformer~\cite{Dosovitskiy2021ViT} --- ingests the request's images as
grids of patches and emits embedding tokens. Computationally this stage
resembles neither prefill nor decode: its attention is non-causal and
carries no KV cache; its cost scales with image count and resolution
rather than conversation length; and its work arrives in bulk --- in the
workload studied here, images arrive as $k\times k$ tile and frame grids
of up to $97^2=\num{9409}$ images in a single request. These properties
invite \emph{modality disaggregation}: run the encoder as its own pool, on
its own hardware class, scaled independently of the language model --- the
architecture demonstrated by ModServe~\cite{Qiu2025ModServe} and now
pursued by fully disaggregated multimodal serving
runtimes~\cite{vLLMOmni2026} and elastic multimodal serving
frameworks~\cite{Liu2025ElasticMM}. This thesis adopts modality disaggregation
as its system model and extends it on the control side.

Independent pools require independent provisioning. The prevailing
practice, inherited from stateless web serving, is \emph{reactive}: watch
an aggregate load signal --- the request rate, or a utilization average
--- and add or remove replicas when a threshold is crossed. Two
assumptions hide inside that loop: that the watched signal actually
carries each pool's demand, and that new replicas arrive faster than the
load changes --- optimistic for GPU serving, where bringing up a replica
(scheduling, weight loading, warmup) is measured in tens of seconds to
minutes~\cite{Fu2024ServerlessLLM}.\footnote{The lead time is
pool-specific. Measured on this thesis's benchmark node, the 576M-parameter
vision encoder reaches serving-ready in ${\approx}$3--5\,s of
\emph{model-side} work --- dominated by interpreter and framework imports,
with the 1.15\,GB checkpoint costing only 0.11\,s to read (O\_DIRECT NVMe)
plus 0.13\,s of host-to-device transfer --- so encoder-pool lead times
above ${\approx}5$\,s measure orchestration (pod scheduling, image pull,
node acquisition), not the model. Language-pool workers, with an order of
magnitude more weights plus engine initialization and KV-cache allocation,
sit in the tens-of-seconds-and-up regime of the systems cited above. The
asymmetry is structural, not model-specific: across current VLM families
the vision tower stays at hundreds of millions to a few billion
parameters while decoders scale to hundreds of billions, so the encoder
pool's lead is orchestration-bound in general --- and the gap widens with
every decoder generation. The
closed-loop evaluation therefore sweeps per-pool lead-time ranges rather
than one shared constant.} Prediction-driven capacity management
has begun to appear for text-only LLM fleets --- forecast-aware fleet
planning~\cite{Jaiswal2025SageServe}, prewarming GPUs against
forecast peak load~\cite{WarmServe2025} --- but it watches aggregate,
request-level signals; what signal a predictor should watch in a
modality-disaggregated system is precisely what is at issue. This thesis
examines both assumptions against production multimodal traffic and finds
that, for the vision pool, both fail.

\section{The provisioning gap}
\label{sec:intro-gap}

The empirical basis of this thesis is
\textit{AzureLMM\-Inference\-Dataset2025}, the production trace Microsoft
released with ModServe~\cite{AzureLMM2025,Qiu2025ModServe}: one
million requests over exactly seven days of a multimodal inference cluster, each row carrying an arrival
timestamp, an image count, and context- and generated-token counts. Its
structure is consequential for serving. Exactly half the requests are
text-only. Image-bearing requests carry $k\times k$ grids, and the
\emph{frames} class of video-like requests --- grids larger than
$10\times10$, only 2.7\% of all requests --- accounts for 81.2\% of all
image volume. The trace's four demand signals --- requests, images,
context tokens, and generated tokens per second --- are the trace-derived
\emph{offered-work} signals available to a scaling policy for this
system, and they map onto the pools directly: the image rate drives the
vision pool, the context-token rate the language pool's prefill stage,
the generated-token rate its decode. (A deployed controller also
consumes system state --- queue depths, utilization, observed latency ---
but those are consequences of the offered work these four signals carry;
they are necessary demand features, not a complete control state.)

Chapter~\ref{cap:trace-eda} establishes that these signals are not
interchangeable. A constant-fraction model --- vision work as a fixed
multiple of request rate --- explains only $R^2=\num{0.123}$ of the
image-rate series, while the same construction explains $R^2=\num{0.654}$
of decode demand (measured as offered work stamped at arrival --- an
upper bound on decode load; Chapter~\ref{cap:trace-eda}): the request
signal carries most of the information about
decode demand and almost none about vision demand. The image signal is
also an order of magnitude burstier (CV$^2$ 3.83 versus 0.32 at one-minute
resolution), and its extremes do not co-occur with the request signal's:
the week's largest burst drives image load to \SI{370}{img/s} --- nearly
sixteen times its weekly mean --- while the request rate in those same
minutes peaks below \num{4.5} requests per second, unremarkable in level
(barely above its own p99 of 4.35); more decisively, over the whole week
the two signals' top-percentile minutes are essentially disjoint (Jaccard
$0.02$). Each major burst is, moreover, a near-pure flood of the frames
class (93--96\% of episode image volume): the vision pool's load is
dominated by a workload class whose arrival the request rate barely
registers. The signal a rate-based autoscaler watches does not carry this
pool's load --- although the load is knowable: every request declares its
image count at admission, so watching the right signal requires no new
instrumentation, only aggregating what the request stream already
contains.

Nor would watching the right signal reactively suffice, because the bursts
are fast and the burstiness does not average away. Burst onsets ramp from
baseline to peak in a median of six minutes --- the fastest in one ---
with a median steepest single-minute jump of \SI{208}{img/s}; against
replica cold starts of tens of seconds to minutes, the fast half of the
bursts cannot be caught by a policy that waits to observe them.
Aggregation does not rescue the reactive loop either: the arrival process
is overdispersed at every operational scale --- the Fano factor of the
counting process grows from 1.8 at one second to 938 at one hour --- so no
averaging window turns the signal into something a threshold rule tracks
safely: the scale-free burstiness long documented for network and Web
traffic~\cite{Leland1994,Crovella1997} reappears in a serving system's
demand signals. The failure is also invisible to the dashboards reactive systems
actually watch: in the two-pool simulation of Chapter~\ref{cap:twopool},
the encoder pool violates its tail-latency target \emph{at the trace's own
recorded rate}, with p99 encoder delay reaching \SI{5.0}{\second} while
mean utilization sits at 29\% --- chunked frames-class requests induce
episodic head-of-line blocking behind the pool's FIFO, and a monitor
watching request rate, mean utilization, or any minute-averaged series
sees a healthy pool all day.

The same analysis identifies the remedy. The signals that carry the
vision pool's load are not only the burstiest but also the most
persistent: the image signal retains autocorrelation
with a 69-minute half-life, against nine minutes for requests, and its
spectral long-range-dependence exponent is the strongest of the four
signals (Hurst ${\approx}0.93$ versus $0.83$ for requests, with
excess dispersion that keeps growing across every aggregation window
measured). The signal most dangerous to reactive scaling is
thus also the most forecastable one, and it is forecastable at exactly the
horizons provisioning needs: minutes to tens of minutes, the scale of
burst onsets and replica cold starts. The provisioning gap is a lead-time
problem, and the lead time is available --- to a forecaster.

\section{Thesis statement and research questions}
\label{sec:intro-thesis}

The position this document defends can be stated compactly:

\begin{quote}
\itshape In a modality-disaggregated multimodal serving system, no single
load signal can size every pool: the pools' demands originate in different
request properties, decouple in time, and saturate their pools at offered
loads far apart. Each pool must therefore be provisioned on its own demand
signal --- and can be provisioned \emph{predictively}, because the signal
least visible to conventional autoscaling is also the most persistent:
pretrained time-series models forecast it, zero-shot, at the horizons
provisioning requires. Scaling on these forecasts holds the tail latency
objective through the burst onsets that defeat reactive control at every
tested cold-start lead time, at a fraction of the GPU-hours of the best
static fleet.
\end{quote}

Four research questions structure the work.

\begin{description}
\item[RQ1 --- Signal sufficiency.] Does the request-rate signal, the input
of conventional autoscalers, carry the information needed to provision
each pool of a modality-disaggregated system?
Chapters~\ref{cap:trace-eda} and~\ref{cap:twopool} answer no for the
vision pool --- statistically ($R^2=\num{0.123}$) and then mechanically
(saturation knees an order of magnitude apart on steady traffic; inside
bursts the vision pool fails invisibly at the trace's own rate) --- and
largely yes for the decode stage ($R^2=\num{0.654}$, on the offered-work
reading of the decode signal).

\item[RQ2 --- Forecastability.] Are the per-pool demand signals
forecastable at provisioning horizons, and do pretrained time-series
foundation models realize that forecastability zero-shot --- with what
accuracy, what degradation across the horizon, and what value from
cross-signal covariates and quantile bands? Chapter~\ref{cap:predictors}
answers yes: three foundation models beat every classical baseline in all
twenty (resolution, signal) cells; accuracy decays only marginally across
a twelve-step horizon, which covers the median burst onset; covariates are
nearly irrelevant at provisioning cadence --- the forecasting-side echo of
RQ1's decoupling --- and the quantile band's role is headroom rather than
point correction.

\item[RQ3 --- Platform validity.] Can the two-pool system be simulated
defensibly enough to carry autoscaling conclusions --- encoder costs from
physically validated silicon data, language-pool timing from a calibrated
engine, and per-pool saturation attributable by construction?
Chapters~\ref{cap:aiconfigurator} and~\ref{cap:twopool} build that
platform: the collected encoder-attention operator agrees with physical
measurement within ${\approx}3\%$, the end-to-end kernel ledger closes
with no unexplained residual, and the composed forward carries an
attributed, conservative-direction measurement band
(${\approx}1.4$--$1.6\times$) rather than an assumed accuracy; the language pool's calibrated timing
exposes the latency tails the simulator's default polynomial flattens
(TTFT p99 $+24\%$ on the steady window; on the burst window the
polynomial misses a $1.9\times$ SLA violation outright).

\item[RQ4 --- Closed-loop payoff.] Does provisioning each pool on a
forecast of its own signal improve SLA compliance and GPU-hours over the
two reference policies --- the best static configuration and a reactive
planner --- during the trace's burst episodes? The completed experiments
(chapters in preparation) answer yes, and locate \emph{why}: on the
fastest-onset burst the forecast-driven encoder schedule holds the tail
bound (p99 \SI{588}{\milli\second}) where a reactive controller watching
the \emph{same} image signal fails at every tested cold-start lead from
30 to \SI{300}{\second} (p99 $\geq$\SI{38}{\second}) --- regime
anticipation, not signal choice, is the operative advantage; on the calm
day the same policy runs at a tenth of the deployable static fleet's
provisioned worker-seconds; and the reactive planner's failure mode on
the language pool is diagnostic --- all fourteen configurations satisfy
their own window-mean control law while violating the tail SLA, so
mean-based capacity conversion, not loop tuning, is what fails.
\end{description}

{\renewcommand{\sectionmark}[1]{}%
\section[Approach: trace-driven simulation, anchored in
silicon]{Approach: trace-driven simulation,\protect\\ anchored in silicon}%
\label{sec:intro-approach}}
\sectionmark{Approach}

The questions above are systems questions, but they cannot be answered
responsibly on a production cluster: the experiments sweep offered load
past saturation, replay the same week under many worker configurations and
control policies, and require counterfactuals --- the same trace, the same
arrivals, a different scaler. The methodology is therefore trace-driven,
discrete-event simulation --- the established instrument for
systems questions at this scale in LLM serving~\cite{Agrawal2024Vidur} ---
with fidelity treated as a measured obligation rather than an assumption: every timing model in the platform is either
collected on silicon and validated against physical measurement, or
compared against its alternative so that its effect on the conclusions is
known.

The platform prices each pool on the hardware class it targets. The
\emph{vision pool} runs as an event-driven FIFO of encoder workers whose
per-forward cost comes from \textit{aiconfigurator} (AIC), NVIDIA's
offline deployment planner, querying performance data collected for this
thesis on an NVIDIA GeForce RTX~5090: the encoder is a 576M-parameter
tower whose forward pass fits comfortably in consumer memory, making
consumer-class parts a plausible, cost-effective hardware class for an
independently scaled encoder pool. No encoder-timing data existed for any
consumer or workstation GPU, and AIC has no usable fallback for it (its
speed-of-light bound is ${\approx}5.9\times$ optimistic, and
support-matrix entries on uncollected systems silently certify text-only
costing), so the two operator families that dominate encoder cost ---
non-causal attention and bf16 GEMM --- were collected at locked clocks
inside the pinned production container and contributed upstream (draft
pull request ai-dynamo/aiconfigurator~\#1312). The \emph{language pool} is
simulated by Dynamo's discrete-event engine (\texttt{dynamo.replay} with a
KV-aware router)~\cite{Dynamo2025} under AIC-calibrated H100 timing ---
the decoder-pool platform this architecture targets. The two pools compose
through the trace itself: each request's encoder delay (queue wait,
service, handoff transfer) is folded into its language-pool arrival
timestamp, a seam that keeps per-pool saturation attributable because
either pool can be run in isolation on the same workload.
Figure~\ref{fig:intro-arch} sketches the system and where each demand
signal acts on it.

\begin{figure}[htbp]
\centering
\resizebox{\linewidth}{!}{%
\begin{tikzpicture}[
  font=\small,
  box/.style={draw=black!55, rounded corners=2pt, align=center,
              inner sep=5pt, fill=black!3},
  pool/.style={box, fill=blue!6, minimum height=1.9cm},
  ctrl/.style={box, fill=orange!8},
  arr/.style={-{Stealth[length=2.2mm]}, thick, black!65},
  lbl/.style={font=\scriptsize, black!60, align=center},
]
\node[box] (trace) {Azure LMM trace\\\scriptsize 1M requests / 7 days\\
  \scriptsize (timestamp, images,\\\scriptsize ctx tokens, gen tokens)};
\node[pool, right=1.6cm of trace] (vis) {\textbf{Vision pool}\\
  \scriptsize $C(t)$ encoder workers (RTX 5090)\\
  \scriptsize AIC-costed $f(N)$, FIFO,\\
  \scriptsize chunk-parallel dispatch\\
  \scriptsize SLO: enc.\ delay p99 $<$ 1 s};
\node[pool, right=2.6cm of vis] (llm) {\textbf{Language pool}\\
  \scriptsize $N \times$ TP H100 workers\\
  \scriptsize \texttt{dynamo.replay}, KV router\\
  \scriptsize prefill + decode\\
  \scriptsize SLO: TTFT p99 $<$ 1 s,\\
  \scriptsize ITL p99 $<$ 100 ms};
\node[ctrl, above=1.25cm of $(vis.north)!0.5!(llm.north)$] (fc)
  {\textbf{Per-pool forecaster} (zero-shot TSFM)\\
  \scriptsize image rate $\to C(t)$ \quad token rates $\to N(t)$\\
  \scriptsize quantile headroom, cold-start lead};
\draw[arr] (trace) -- node[lbl, above] {requests} (vis);
\draw[arr] (vis) -- node[lbl, above] {seam:\\arrival $+$ delay}
  node[lbl, below] {$196{\times}4096$\\bf16/img} (llm);
\draw[arr] (fc.west) -| node[lbl, left, pos=0.85] {$C(t)$} (vis.north);
\draw[arr] (fc.east) -| node[lbl, right, pos=0.85] {$N(t)$} (llm.north);
\draw[arr, dashed] (trace.north) |- node[lbl, above, pos=0.72]
  {demand signals:\\req/s, img/s, ctx tok/s, gen tok/s} (fc.west);
\end{tikzpicture}}
\caption{The two-pool system under study. The trace drives an
AIC-costed encoder pool composed with Dynamo's discrete-event language
pool through a timestamp-shift seam; a per-pool forecaster converts each
pool's own demand signal into its worker schedule under cold-start lead
times. Solid arrows carry requests; the dashed arrow carries the
materialized demand signals the controller consumes.}
\label{fig:intro-arch}
\end{figure}

The model under study is the Qwen3-VL family~\cite{Qwen3VL2025}, chosen
for a property that
makes encoder-side results unusually portable: its 8B, 32B, and 235B
variants share one vision tower, and the 8B and 235B also share its
4096-dimensional projector, so the calibrated encoder
model, its batching law, and the burst-saturation result transfer across
the family (for the 32B, whose projector targets its 5120-dimensional
embedding space, the final projector GEMM and the handoff bytes scale by
$5/4$ while the tower-dominated cost is unchanged); a spot-check prices the
235B flagship at
tensor-parallel~8 on the same shipped H100 data and finds its saturation
knee one swept octave below the 8B's --- decoder scale halves the language
pool's operating envelope while leaving the vision pool untouched. Forecasting uses three
pretrained time-series foundation models ---
TimesFM-2.0~\cite{Das2024TimesFM}, Chronos-2~\cite{Ansari2025Chronos2},
and Toto-2.0~\cite{Khwaja2026Toto2} --- against naive,
exponential-smoothing, and Holt--Winters baselines, under a single
rolling-origin protocol whose stitched forecasts tile the timeline exactly
as a scaler would consume them.

One methodological commitment runs through all of it: every number in this
document is generated by a versioned script over archived inputs, every
collection and validation phase leaves an audit artifact in the workspace,
the collected datasets were offered upstream, and corrections made during
the work are recorded in the affected chapters rather than silently
absorbed.

\section{Contributions}
\label{sec:intro-contrib}

\begin{enumerate}
\item \textbf{A characterization of production multimodal inference
traffic as per-pool demand signals} (Chapter~\ref{cap:trace-eda}) --- to
our knowledge the first such view of a production multimodal trace,
possible only since the trace's 2025 release with ModServe.
GPU-accelerated, fully scripted analysis establishing:
the decoupling of vision work from request rate ($R^2=\num{0.123}$, versus
\num{0.654} for decode); the burst anatomy (single-class frames-floods,
93--96\% of episode volume; onsets with a median of six minutes and
single-minute jumps of \SI{208}{img/s}); the persistence contrast (Hurst
${\approx}0.83$ versus $0.93$; autocorrelation half-lives of 9 versus 69
minutes for requests versus images); and the workload facts any converter must honor (the $k^2$ grid
structure, the 50\% text-only split, and vision tokens embedded
non-constantly in the context-token field).

\item \textbf{A zero-shot benchmark of time-series foundation models on
serving demand signals} (Chapter~\ref{cap:predictors}) --- to our
knowledge the first evaluation of these models framed as an autoscaler
would consume them. TimesFM, Chronos-2, and Toto-2 against three
classical baselines over all four signals at five resolutions under one
rolling-origin protocol: the
foundation models win all twenty cells; twelve steps of lead time cost
only ${\approx}2$ WAPE points at one-minute resolution; covariates are
nearly irrelevant at provisioning cadence; no fixed quantile improves
on a median point forecast, leaving the band as headroom --- and the p90
band edge is empirically calibrated at provisioning cadence (coverage
0.88--0.90 across all models and signals). The full forecast set ---
21.4 million forecast rows, each a point, nine quantiles, and the
realized value --- is archived alongside the metric tables.

\item \textbf{A silicon-calibrated cost model for consumer-GPU vision
encoding} (Chapter~\ref{cap:aiconfigurator}) --- the planner's first
encoder performance dataset for any non-datacenter GPU, a gap its own
query modes turn into hard errors. The RTX~5090 onboarded as a new AIC
system; encoder-attention and bf16-GEMM datasets collected at locked
clocks in the production container and upstreamed (draft PR
ai-dynamo/aiconfigurator~\#1312); a validation methodology that closes the
kernel-time ledger against a physical measurement of the same network on
the same silicon (attention within ${\approx}3\%$; every residual term
attributed; the composed forward bracketed by a conservative-direction
measurement band);
and two findings of independent value to the tool's ecosystem --- AIC has
no usable fallback for encoder timing, so support-matrix ``PASS'' entries
on uncollected systems certify text-only costing, and unlocked
consumer-GPU clocks silently inflate collected data by ${\approx}17\%$
while leaving validation ratios unchanged.

\item \textbf{A calibrated two-pool simulation platform, and the
saturation-knee result} (Chapter~\ref{cap:twopool}). Composition of the
AIC-costed encoder pool with Dynamo's discrete-event language pool through
a timestamp-shift seam with per-pool attribution; a probe-verified exact
batching law for the encoder (cost depends only on the total images per
forward) with an analytic inter-pool handoff bound; quantification of the
calibrated-versus-polynomial timing delta (steady TTFT p99 $+24\%$; a
burst-window SLA violation the polynomial misses outright); and the
central result: the pools' saturation knees sit an order of magnitude
apart on steady traffic, and the largest burst breaks the two pools in
structurally different ways --- the language pool through a
shape-fixable prefill floor, the encoder pool through episodic queueing
at the trace's own rate and 29\% mean utilization, invisible to rate-
and utilization-based monitors and fixable only by capacity. The
decoupling of Chapter~\ref{cap:trace-eda}, reproduced as system behavior.

\item \textbf{Baselines and the closed predictive-scaling loop}
(chapters in preparation; experiments complete, archived, and
summarized here). The static-optimal worker-grid frontier per window
(the burst windows need 6--16 encoder workers at trace rate against 2
on a calm day, and TP2 language shapes throughout); a reactive
SLA-planner baseline whose fourteen language-pool configurations all
satisfy their own window-mean control law while violating the tail SLA
--- isolating mean-based capacity conversion as the defect; and the
contribution proper: per-pool forecasts driving worker schedules under
cold-start lead times, scored as SLA compliance $\times$
provisioned worker-seconds. On the fastest-onset burst the forecast
schedule holds the tail bound (p99 \SI{588}{\milli\second}) where the
reactive controller --- watching the \emph{same} per-pool signal ---
fails at every lead time from 30 to \SI{300}{\second}, degrading
monotonically toward its never-scaled limit (p99 $\geq$\SI{38}{\second});
on the calm day the forecast schedule spends a tenth of the deployable
static fleet's provisioned worker-seconds. Regime anticipation ---
raising capacity before an onset a trailing window cannot see --- is
identified as the operative mechanism, with quantile choice (p90) as
the headroom dial.
\end{enumerate}

\section{Scope and limitations}
\label{sec:intro-scope}

The claims are scoped by their evidence. The workload is one seven-day
production trace with a by-construction 50/50 modality split and
arrival-side information only: the rate signals stamp each request's full
token volume at its arrival minute, which is faithful for encoder and
prefill work but makes the decode series an offered-work upper bound. The
trace records image counts, not sizes, so the encoder pool assumes a
uniform $448^2$ image, with a measured, bounded conservative bias
(${\approx}1.14\times$ on comparable work) documented where it matters.
The simulation is open-loop and deterministic; the inter-pool seam
serializes each request's encoder and prefill phases (no speculative
overlap); the latency targets used for knee detection are the thesis's evaluation
SLA (TTFT and encoder-delay p99 under \SI{1}{\second}, inter-token p99
under \SI{100}{\milli\second}; Chapter~\ref{cap:twopool}), a design choice
whose sensitivity is quantified there; and language-pool results are reported for the 8B model at
TP1 with a 235B-TP8 spot-check, leaving intermediate family members
unpriced. Forecasts are zero-shot --- no fine-tuning on the trace ---
which is the conservative reading of the foundation models' capability.
\emph{[Placeholder: the planned validation anchors --- the winning policy
replayed against the live Dynamo runtime with mocked workers, and
optionally one physical encoder worker on the RTX~5090 under a burst
replay --- and what they did or did not confirm.]}

\section{Organization}
\label{sec:intro-org}

Chapter~\ref{cap:related} situates the thesis in the five literatures it
borders --- LLM serving, multimodal serving, autoscaling, workload
characterization, and time-series forecasting --- and locates the empty
intersection the contribution occupies. The rest of the document proceeds
in two movements. The first is empirical.
Chapter~\ref{cap:trace-eda} characterizes the trace and establishes the
decoupling, the burst anatomy, and the persistence contrast that make
per-pool forecasting both necessary and plausible.
Chapter~\ref{cap:predictors} tests the forecasting premise directly,
benchmarking three time-series foundation models against classical
baselines on the materialized demand signals. The second movement builds
and exercises the system. Chapter~\ref{cap:aiconfigurator} constructs and
physically validates the vision pool's cost model on collected RTX~5090
silicon. Chapter~\ref{cap:twopool} composes the two-pool simulation,
quantifies both pools' timing fidelity, and locates their saturation knees
--- the decoupling as measured system behavior. Four chapters are in
preparation from completed and archived experiments: a \emph{baselines}
chapter establishing the static-optimal worker grid and the reactive
SLA-planner baseline on both pools; a \emph{closed-loop} chapter scoring
forecast-driven per-pool schedules against those baselines as SLA
compliance $\times$ provisioned worker-seconds, with ablations over
cold-start lead, headroom quantile, and forecasting model; a
\emph{validation} chapter replaying the winning policy against the live
Dynamo runtime with mocked workers; and \emph{conclusions}.
